\documentclass{../assignment}

\عنوان{عملی \متن‌لاتین{HTTP}}

\ترم{پاییز ۱۴۰۵}
\مهلت{۱۱ آبان ۱۴۰۵}
\تله‌هوش‌مصنوعی

\begin{document}

\عنوان‌ساز

\فهرست‌مطالب

\قسمت{مقدمه}

هدف این تمرین پیاده‌سازی یک \متن‌لاتین{reverse proxy} با احراز هویت \متن‌لاتین{JWT} برای یک سرویس قدیمی است.
مجموع نمرهٔ تمرین ۲۰ است.

\قسمت{عملی}

\زیرقسمت{احراز هویت با پراکسی}

از شما خواسته شده تا یک \متن‌لاتین{reverse proxy} برای شبکه داخلی یک سرویس قدیمی \متن‌لاتین{(legacy)} طراحی کنید.
سرویس قدیمی امکان \متن‌لاتین{authentication} نداشته و این مورد برای استفاده اعضای داخل سازمان از شبکه داخلی مشکلی ندارد، ولی انتشار این سرویس بر روی اینترنت یک ضعف امنیتی بزرگ است چرا که همه امکان استفاده از این سرویس را دارند.

پروکسی خود را بر سر ورودی شبکه اینترنت قرار داده و آن را به گونه‌ای تنظیم کنید که همه درخواست‌های پورت ۸۰ را از شبکه خارجی دریافت کرده و به یک پورت (مثلا ۸۰۸۱) در شبکه داخلی ارسال کند.

برای احراز هویت درخواست ها در پروکسی حتما از \متن‌لاتین{\متن‌سیاه{JWT authentication}} استفاده کنید.

برای حل این تمرین به موارد زیر پاسخ دهید:

\شروع{شمارش}
\فقره چگونه می‌توان یک سرویس را فقط بر روی شبکه داخلی اجرا کرد؟ \نمره{۵}
\فقره علاوه بر \متن‌لاتین{jwt authentication} یکی دیگر از راه های امن کردن سرویس داخلی \متن‌لاتین{legacy}، استفاده از روشی به نام \متن‌لاتین{TLS Termination} است. فلسفه این کار را توضیح دهید. \نمره{۵}
\فقره یک \متن‌لاتین{jwt authentication proxy} آماده کنید و برای \تارنما{https://github.com/mccutchen/go-httpbin}{این پروژه} از آن استفاده کنید. این پروژه نسخهٔ \متن‌لاتین{Go} و به‌روز \متن‌لاتین{httpbin} است و ایمیج آن با نام \متن‌لاتین{\texttt{ghcr.io/mccutchen/go-httpbin}} در دسترس است. \نمره{۱۰}
\پایان{شمارش}

این سرویس می‌تواند توسط خود دانشجو برنامه‌نویسی شود و یا توسط پیکربندی وب سرورها و \متن‌لاتین{reverse proxy}های آماده صورت بپذیرد.

\قسمت{نحوهٔ تحویل}

پاسخ دو مورد اول را در فایل \متن‌لاتین{README.md} مخزن خود، با شمارهٔ بخش مربوط، بنویسید. برای بخش سوم کد و یا پیکربندی‌های موجود در مخزن بررسی خواهد شد.
دقت بفرمایید که فایل \متن‌لاتین{PDF} قابل قبول نمی‌باشد.

\قوانین

\پایان‌ساز

\end{document}
